% Part: turing-machines
% Chapter: machines-computations
% Section: representing-tms

\documentclass[../../../include/open-logic-section]{subfiles}

\begin{document}

\olfileid{tur}{mac}{rep}
\olsection{ट्यूरिंग यंत्रांचे निरूपण}

\begin{explain}
ट्यूरिंग यंत्रे \emph{अवस्था-आलेखांद्वारे} चित्ररूपात दाखवता येतात.
या आलेखांत बाणांनी जोडलेली अवस्थांची घरे असतात. अर्थातच प्रत्येक
अवस्था-घर यंत्राची एक अवस्था दर्शवते. प्रत्येक बाण त्या अवस्थेतून
अमलात आणता येणारी एक सूचना दर्शवतो; सूचनेचे तपशील त्या बाणाच्या
वर किंवा खाली लिहिलेले असतात. पुढील यंत्र पाहा. त्याच्या $q_0$ आणि
$q_1$ अशा फक्त दोन अंतर्गत अवस्था आणि एकच सूचना आहे:
\[
\begin{tikzpicture}[->,>=stealth',shorten >=1pt,auto,node distance=2.8cm,
                    semithick]
  \tikzstyle{every state}=[fill=none,draw=black,text=black]

  \node[initial,state]   (A)              {$q_0$};
  \node[state]   (B) [right of=A] {$q_1$};

  \path (A) edge  node {\TMtrans{\TMblank}{\TMstroke}{\TMright}} (B);
\end{tikzpicture}
\]
ट्यूरिंग यंत्राला वाचन-लेखन शीर्ष आणि त्यावर आदान लिहिलेली फीत
असते, हे आठवा. या सूचनेचा अर्थ असा: \emph{$q_0$ अवस्थेत
रिकामे चिन्ह~$\TMblank$ वाचत असाल, तर $\TMstroke$ लिहा,
उजवीकडे सरका आणि $q_1$ अवस्थेत जा}. हे संक्रमण फलनाने
$\tuple{q_0, \TMblank}$ याचा $\tuple{q_1, \TMstroke,
\TMright}$ शी केलेला संबंध याच्या सममूल्य आहे.
\end{explain}

\begin{ex}
\emph{सम यंत्र}: फितीवर $\TMstroke$ चिन्हांची संख्या सम असेल तेव्हा
आणि तेव्हाच पुढील ट्यूरिंग यंत्र थांबते (फितीवरील पहिल्या
$\TMblank$ आधी सर्व $\TMstroke$ चिन्हे येतात, हे गृहीत धरून).
\[
\begin{tikzpicture}[->,>=stealth',shorten >=1pt,auto,node distance=2.8cm,
                    semithick]
  \tikzstyle{every state}=[fill=none,draw=black,text=black]

  \node[initial,state]         (A)              {$q_0$};
  \node[state]         (B) [right of=A] {$q_1$};

  \path (A) edge [bend left] node {\TMtrans{\TMstroke}{\TMstroke}{\TMright}} (B)
        (B) edge [loop above] node {\TMtrans{\TMblank}{\TMblank}{\TMright}} (B)
            edge [bend left] node {\TMtrans{\TMstroke}{\TMstroke}{\TMright}} (A);
\end{tikzpicture}
\]

हा अवस्था-आलेख पुढील संक्रमण फलनाशी संबंधित आहे:
\begin{align*}
\delta(q_0, \TMstroke) & = \tuple{q_1, \TMstroke, \TMright},\\
\delta(q_1, \TMstroke) & = \tuple{q_0, \TMstroke, \TMright},\\
\delta(q_1, \TMblank)  & = \tuple{q_1, \TMblank, \TMright}
\end{align*}
\end{ex}

\begin{explain}
वरील यंत्र आदानातील रेघांची संख्या सम असेल तेव्हाच थांबते.
अन्यथा यंत्र (सैद्धांतिकदृष्ट्या) अमर्याद काळ कार्यरत राहते.
कोणत्याही यंत्राच्या कोणत्याही आदानासाठी त्याच्या \emph{स्थितिवर्णनांचा}
क्रम पाहून प्रदान निश्चित करता येते. स्थितिवर्णनाची औपचारिक
व्याख्या पुढे देऊ. सध्या, यंत्र चालू असताना प्रत्येक क्षणी त्याची
स्थिती दाखवणाऱ्या आकृत्यांची मालिका असे स्थितिवर्णनांचे अंतर्ज्ञानी
आकलन करता येईल. स्थितिवर्णनात फितीवरील मजकूर, यंत्राची अवस्था
आणि वाचन-लेखन शीर्षाची जागा दिसते.

सम यंत्राला चार $\TMstroke$ चिन्हांचे आदान दिल्यास त्याची
स्थितिवर्णने पाहू. या वेळी यंत्र थांबेल, अशी अपेक्षा आहे. मग
तीन $\TMstroke$ चिन्हांचे आदान देऊ; त्या वेळी ते कायम चालू राहील.

यंत्र $q_0$ अवस्थेत सुरू होते आणि सर्वांत डावीकडील $\TMstroke$
वाचते. त्याची प्रारंभिक स्थिती अशी दर्शवता येते:
\[
\TMendtape \TMstroke_0 \TMstroke \TMstroke \TMstroke \TMblank \ldots
\]
वरील स्थितिवर्णन सरळ आहे. दिसते त्याप्रमाणे यंत्र $q_0$ अवस्थेत
सुरू होते आणि सर्वांत डावीकडील $\TMstroke$ वाचते. पहिल्या
%READERNOTE{OLTUR-001 — गोठवलेल्या इंग्रजीत आरंभीची अवस्था one म्हटली आहे; लगेच आधीचे चित्र आणि संक्रमण फलन शून्य अवस्था दाखवतात. मराठीत त्यांच्याशी सुसंगत शून्य अवस्था वापरली आहे.}
$\TMstroke$ च्या खाली लिहिलेला अवस्थेच्या नावाचा निर्देशांक हे
दर्शवतो. या टप्प्यावर लागू होणारी सूचना
$\delta(q_0, \TMstroke) = \tuple{q_1, \TMstroke, \TMright}$ आहे;
म्हणून यंत्र फितीवर उजवीकडे सरकते आणि $q_1$ अवस्थेत जाते.
\[
\TMendtape \TMstroke \TMstroke_1 \TMstroke \TMstroke \TMblank \ldots
\]
आता यंत्र $q_1$ अवस्थेत $\TMstroke$ वाचत आहे, म्हणून
$\delta(q_1, \TMstroke) = \tuple{q_0,
  \TMstroke, \TMright}$ ही सूचना ``पाळावी'' लागते. त्यातून पुढील
स्थितिवर्णन मिळते:
\[
\TMendtape \TMstroke \TMstroke \TMstroke_0 \TMstroke \TMblank \ldots
\]
यंत्र पुढे चालू राहिल्यावर हे नियम त्याच क्रमाने पुन्हा लागू होतात
आणि पुढील दोन स्थितिवर्णने मिळतात:
\[
\TMendtape \TMstroke \TMstroke \TMstroke \TMstroke_1 \TMblank \ldots
\]
\[
\TMendtape \TMstroke \TMstroke \TMstroke \TMstroke \TMblank_0 \ldots
\]
आता यंत्र $q_0$ अवस्थेत $\TMblank$ वाचत आहे. संक्रमण-आलेखावरून
येथे अमलात आणण्यासाठी कोणतीही सूचना नाही, हे सहज दिसते;
म्हणून यंत्र थांबले आहे. याचा अर्थ आदान स्वीकारले गेले.

आता यंत्राला तीन $\TMstroke$ चिन्हांचे आदान देऊन सुरू करू.
सुरुवातीची काही स्थितिवर्णने मिळतीजुळती असतील, कारण त्याच सूचना
अमलात येतात; फक्त फितीवरील आदान थोडे वेगळे आहे:
\[
\TMendtape \TMstroke_0 \TMstroke \TMstroke \TMblank \ldots
\]
\[
\TMendtape \TMstroke \TMstroke_1 \TMstroke \TMblank \ldots
\]
\[
\TMendtape \TMstroke \TMstroke \TMstroke_0 \TMblank \ldots
\]
\[
\TMendtape \TMstroke \TMstroke \TMstroke \TMblank_1 \ldots
\]
आता यंत्र सर्व $\TMstroke$ चिन्हांच्या पुढे गेले आहे आणि $q_1$
अवस्थेत $\TMblank$ वाचत आहे. आलेखात दाखवल्याप्रमाणे
$\delta(q_1, \TMblank) =\tuple{q_1,
\TMblank, \TMright}$ या रूपाची सूचना आहे. फीत उजवीकडे अनंतपर्यंत
$\TMblank$ चिन्हांनी भरलेली असल्याने यंत्र ही सूचना \emph{कायम}
अमलात आणत राहील; ते $q_1$ अवस्थेतच राहून सतत उजवीकडे सरकेल.
यंत्र कधीच थांबणार नाही आणि आदान स्वीकारणार नाही.
\end{explain}

\begin{explain}
सर्व यंत्रे थांबतातच असे नाही, हे लक्षात ठेवणे महत्त्वाचे आहे. थांबणे
म्हणजे अमलात आणण्यासाठी सूचना संपणे असेल, तर प्रत्येक अवस्थेत
प्रत्येक चिन्हासाठी बाहेर जाणारा बाण ठेवून कधीही न थांबणारे यंत्र
बनवता येते. $q_0$ अवस्थेत $\TMblank$ वाचण्यासाठी एक सूचना
जोडून सम यंत्र कायम चालू राहील असे बदलता येते.
\end{explain}

\begin{ex}
\[
\begin{tikzpicture}[->,>=stealth',shorten >=1pt,auto,node distance=2.8cm,
                    semithick]
  \tikzstyle{every state}=[fill=none,draw=black,text=black]

  \node[initial,state] (A)              {$q_0$};
  \node[state]         (B) [right of=A] {$q_1$};

  \path
  (A) edge [bend left] node {\TMtrans{\TMstroke}{\TMstroke}{\TMright}} (B)
      edge [loop above] node {\TMtrans{\TMblank}{\TMblank}{\TMright}} (A)
  (B) edge [loop above] node {\TMtrans{\TMblank}{\TMblank}{\TMright}} (B)
      edge [bend left] node {\TMtrans{\TMstroke}{\TMstroke}{\TMright}} (A);
\end{tikzpicture}
\]
\end{ex}

\begin{explain}
यंत्र-सारणी हा ट्यूरिंग यंत्रांचे निरूपण करण्याचा दुसरा मार्ग आहे.
यंत्र-सारणीत फितीवरील वर्णमाला $x$-अक्षावर आणि यंत्राच्या अवस्था
$y$-अक्षावर दाखवलेल्या असतात. सारणीतील प्रत्येक अवस्था आणि चिन्ह
यांच्या छेदावर सूचनेचा उरलेला भाग---नवी अवस्था, नवे चिन्ह आणि
हालचालीची दिशा---लिहिलेला असतो. यंत्र कोणत्या अवस्थेत आणि कोणते
चिन्ह वाचताना थांबते हे यंत्र-सारणीवरून सहज ठरवता येते. सारणीत
जिथे नोंद नाही, तिथे यंत्र थांबू शकते. अवस्था-आलेख किंवा सूचनासंच
यांत यंत्र नेमके कुठे थांबेल हे लगेच स्पष्ट होत नाही; पण
यंत्र-सारणीत रिकामी घरे शोधून अशी ठिकाणे पटकन ओळखता येतात.
\end{explain}

\begin{ex}
सम यंत्राची यंत्र-सारणी अशी आहे:
\[
\centering
\begin{tabular}{lllll}
\cline{2-4}
\multicolumn{1}{l|}{}      & \multicolumn{1}{c|}{$\TMblank$}                
& \multicolumn{1}{c|}{$\TMstroke$}     & \multicolumn{1}{c|}{$\TMendtape$}   &  \\ \cline{1-4}
\multicolumn{1}{|l|}{$q_0$} & \multicolumn{1}{l|}{}                          
& \multicolumn{1}{l|}{\TMtrans{\TMstroke}{q_1}{\TMright}} & 
\multicolumn{1}{l|}{\phantom{\TMtrans{\TMstroke}{q_1}{\TMright}}} &  \\ \cline{1-4}
\multicolumn{1}{|l|}{$q_1$} & \multicolumn{1}{l|}{\TMtrans{\TMblank}{q_1}{\TMright}} 
& \multicolumn{1}{l|}{\TMtrans{\TMstroke}{q_0}{\TMright}} & 
\multicolumn{1}{l|}{\phantom{\TMtrans{\TMstroke}{q_1}{\TMright}}} &  \\ \cline{1-4}
\end{tabular}
\]
दिसते त्याप्रमाणे यंत्र $q_0$ अवस्थेत $\TMblank$ वाचताना थांबते.
\end{ex}

\begin{explain}
आतापर्यंत आपण फक्त आदान वाचून ते स्वीकारणारी यंत्रे पाहिली.
परंतु ट्यूरिंग यंत्रे वाचूही शकतात आणि लिहूही शकतात. अशा यंत्राचे
एक उदाहरण (अशी अनेक उदाहरणे असली तरी) म्हणजे \emph{दुप्पट करणारे
यंत्र}. हे यंत्र फितीवरील $n$ $\TMstroke$ चिन्हांचा सलग गट
आदान म्हणून घेते आणि $2n$ $\TMstroke$ चिन्हांचा सलग गट प्रदान करते.
\end{explain}

\begin{ex}
  \ollabel{ex:doubler} दुप्पट करणारे यंत्र बनवण्यापूर्वी ही समस्या
  सोडवण्यासाठी \emph{धोरण} ठरवणे महत्त्वाचे आहे. यंत्राला (आपल्या
  व्याख्येनुसार) त्याने किती $\TMstroke$ चिन्हे वाचली हे लक्षात
  ठेवता येत नाही; म्हणून फितीवरील सर्व $\TMstroke$ चिन्हांचा मागोवा
  घेण्याचा मार्ग हवा. एक मार्ग म्हणजे आदान आणि प्रदान यांमध्ये
  $\TMblank$ ठेवून त्यांना वेगळे करणे. मग यंत्र आदानातील पहिले
  $\TMstroke$ पुसू शकते, आदानाचा उरलेला भाग पार करू शकते,
  $\TMblank$ सोडू शकते आणि दोन नवी $\TMstroke$ चिन्हे लिहू शकते.
  नंतर ते आदानातील दुसरे $\TMstroke$ शोधून त्याच्यासाठीही असेच
  दुपटीकरण करील. आदानातील प्रत्येक एका $\TMstroke$ साठी ते
  प्रदानात दोन $\TMstroke$ चिन्हे लिहील. पुढे जाताना आदान
  पुसल्यामुळे कोणतेही $\TMstroke$ चिन्ह सुटणार नाही किंवा त्याची
  दुप्पट मोजणी होणार नाही, याची खात्री होते. संपूर्ण आदान पुसल्यावर
  फितीवर $2n$ $\TMstroke$ चिन्हे उरतील. परिणामी ट्यूरिंग यंत्राचा
  अवस्था-आलेख \olref{fig:doubler} मध्ये दाखवला आहे.
  \begin{figure}
\[
\begin{tikzpicture}[->,>=stealth',shorten >=1pt,auto,node distance=2.8cm,
                    semithick]
  \tikzstyle{every state}=[fill=none,draw=black,text=black]

  \node[initial,state] (1)              {$q_0$};
  \node[state]         (2) [right of=1] {$q_1$};
  \node[state]         (3) [right of=2] {$q_2$};
  \node[state]         (4) [below of=3] {$q_3$};
  \node[state]         (5) [left of=4]  {$q_4$};
  \node[state]         (6) [left of=5]  {$q_5$};

  \path (1) edge node {\TMtrans{\TMstroke}{\TMblank}{\TMright}} (2)
    (2) edge [loop above] node {\TMtrans{\TMstroke}{\TMstroke}{\TMright}} (2)
      edge node {\TMtrans{\TMblank}{\TMblank}{\TMright}} (3)
    (3) edge [loop above] node {\TMtrans{\TMstroke}{\TMstroke}{\TMright}} (3)
        edge node {\TMtrans{\TMblank}{\TMstroke}{\TMright}} (4)
    (4) edge [loop below] node {\TMtrans{\TMblank}{\TMstroke}{\TMleft}} (4)
        edge node {\TMtrans{\TMstroke}{\TMstroke}{\TMleft}} (5)
    (5) edge [loop below]  node {\TMtrans{\TMstroke}{\TMstroke}{\TMleft}} (5)
        edge              node {\TMtrans{\TMblank}{\TMblank}{\TMleft}} (6)
    (6) edge [loop below] node {\TMtrans{\TMstroke}{\TMstroke}{\TMleft}} (6)
        edge              node {\TMtrans{\TMblank}{\TMblank}{\TMright}} (1);
\end{tikzpicture}
\]
\caption{दुप्पट करणारे यंत्र}
\ollabel{fig:doubler}
\end{figure}
\end{ex}

\begin{prob}
एखादे आदान निवडा आणि \olref[tur][mac][rep]{ex:doubler} मधील दुप्पट
करणाऱ्या यंत्राची स्थितिवर्णने क्रमाने शोधा.
\end{prob}

\begin{prob}
$\{\TMendtape,\TMblank, A, B\}$ ही वर्णमाला असलेले ट्यूरिंग यंत्र
रचा. $A$ आणि $B$ यांच्या कोणत्याही चिन्हमालेत $A$ ची संख्या $B$
च्या संख्येइतकी \emph{आणि} सर्व $A$ सर्व $B$ च्या आधी असतील,
तर यंत्राने ती चिन्हमाला स्वीकारावी, म्हणजे तिच्यावर थांबावे.
$A$ आणि $B$ यांच्या संख्या असमान असतील किंवा सर्व $A$ सर्व $B$
च्या आधी नसतील, तर यंत्राने ती नाकारावी, म्हणजे तिच्यावर थांबू नये.
(उदाहरणार्थ, यंत्राने $AABB$ आणि $AAABBB$ स्वीकाराव्यात; पण
$AAB$ आणि $AABBAABB$ या दोन्ही नाकाराव्यात.)
\end{prob}

\begin{prob}
$\{\TMendtape,\TMblank, A, B\}$ ही वर्णमाला असलेले ट्यूरिंग यंत्र
रचा. त्याने $A$ आणि $B$ यांची कोणतीही चिन्हमाला $\alpha$ आदान
म्हणून घेऊन तिची प्रत तयार करावी आणि $\alpha\alpha$ या रूपाचे
प्रदान द्यावे. (उदाहरणार्थ, $ABBA$ या आदानापासून $ABBAABBA$
हे प्रदान मिळावे.)
\end{prob}

\begin{prob}
\emph{वर्णक्रमात आहे का?}: $\{\TMendtape,\TMblank, A, B\}$ ही
वर्णमाला असलेले ट्यूरिंग यंत्र रचा. $A$ आणि $B$ यांची सांत
क्रमिका आदान दिल्यावर सर्व $A$ सर्व $B$ यांच्या डावीकडे आहेत की
नाही, हे त्याने तपासावे. यंत्राने आदान-चिन्हमाला फितीवर ठेवावी;
ती ``वर्णक्रमात'' असेल तर थांबावे आणि नसेल तर कायम पुनरावर्तनात
चालू राहावे.
\end{prob}

\begin{prob}
\emph{वर्णक्रम लावणारे यंत्र}: $\{\TMendtape,\TMblank, A, B\}$
ही वर्णमाला असलेले ट्यूरिंग यंत्र रचा. $A$ आणि $B$ यांची सांत
क्रमिका आदान म्हणून घेऊन त्यांची पुनर्रचना अशा रीतीने करावी की
सर्व $A$ सर्व $B$ यांच्या डावीकडे येतील. (उदाहरणार्थ, $BABAA$
ही क्रमिका $AAABB$ व्हावी आणि $ABBABB$ ही क्रमिका $AABBBB$
व्हावी.)
\end{prob}

\end{document}
